\chapter{Criterio de publicación integral sobria del sistema}

\section{Publicación integral sobria}
\begin{definition}
Se llama publicación integral sobria del sistema al modo disciplinado mediante el cual el corpus MADELEINE puede proyectarse externamente como conjunto sin perder su jerarquía interna, su trazabilidad metodológica ni la diferencia entre fundamento, canon intermedio, capa metodológica y material todavía reservado.
\end{definition}

\begin{observation}
La publicación integral sobria no equivale a apertura indiscriminada ni a exposición total inmediata de todo el material acumulado. Su función es fijar una forma externa del sistema que preserve estructura, límites y sobriedad.
\end{observation}

\section{Condición mínima de publicación integral}
\begin{definition}
Existe publicación integral sobria cuando pueden declararse simultáneamente:
\begin{enumerate}
  \item una arquitectura legible entre Tomo I, Tomo II y Tomo III;
  \item una distinción suficiente entre núcleo, canon operativo, metodología, claims públicos y reservas;
  \item una capa pública mínima capaz de orientarse por fichas, dossiers o mapa externo;
  \item continuidad visible con la genealogía metodológica interna;
  \item ausencia de compresión abusiva del sistema completo en una narrativa externa única.
\end{enumerate}
\end{definition}

\begin{observation}
No se exige todavía mostrar todo el espesor del corpus en una sola salida. Basta con una forma integral sobria que permita comprender la arquitectura general sin falsearla.
\end{observation}

\section{Publicación integral y tomos}
\begin{definition}
Toda publicación integral sobria debe conservar la diferencia funcional entre:
\begin{enumerate}
  \item Tomo I como fijación del núcleo y del canon basal;
  \item Tomo II como canon intermedio y régimen de organización operativa;
  \item Tomo III como capa metodológica, curatorial y de proyección disciplinada de claims.
\end{enumerate}
\end{definition}

\begin{observation}
Esta cláusula impide que la salida externa borre la arquitectura tripartita del sistema o convierta todos los materiales en una sola masa indiferenciada.
\end{observation}

\section{Publicación integral y capa pública}
\begin{definition}
La publicación integral sobria puede apoyarse en:
\begin{enumerate}
  \item fichas públicas mínimas;
  \item dossiers externos de claims seleccionados;
  \item mapa público mínimo del corpus;
  \item referencias explícitas a material observado o reservado cuando resulte necesario para mantener honestidad metodológica.
\end{enumerate}
\end{definition}

\begin{observation}
La capa pública no debe fingir pureza total allí donde el propio proceso metodológico ha preservado observaciones, reservas o gradientes de estabilidad.
\end{observation}

\section{Condición mínima de no publicación integral}
\begin{definition}
No debe intentarse publicación integral sobria cuando la arquitectura externa dependa de simplificaciones que destruyan la diferencia entre fundamento, canon, metodología, capa pública y reserva interna.
\end{definition}

\begin{observation}
La no publicación integral no implica fracaso del sistema. Puede significar solo que todavía conviene fortalecer mediaciones externas más parciales antes de presentar el conjunto como totalidad.
\end{observation}

\section{Cláusula de alcance}
\begin{observation}
El criterio de publicación integral sobria del sistema fija la primera condición de salida externa del conjunto MADELEINE como arquitectura coherente. No constituye todavía una política editorial final exhaustiva ni una teoría cerrada de recepción pública del corpus.
\end{observation}

